%%%%%%%%%%%%%%%%%%%%%%
%%%%%%%%%%%%%%%%%%%%%%
\section{Methods}
%%%%%%%%%%%%%%%%%%%%%%
%%%%%%%%%%%%%%%%%%%%%%

%%%%%%%%%%%%%%%%%%%%%%
%%%%%%%%%%%%%%%%%%%%%%
\subsection{Problem formulation}
%%%%%%%%%%%%%%%%%%%%%%
%%%%%%%%%%%%%%%%%%%%%%

We consider a set of protein sequences $S = \{s_{1}, s_{2}\dots \}$, a set of protein molecular function terms $F = \{f_{1}, f_{2},\dots \}$ and a set of biological process terms $P = \{p_{1}, p_{2}, \dots \}$, where each sequence in $S$ is annotated by at least one element from $F$ or $P$. Additionally, each function $f\in F$ and process $p\in P$ is associated with at least one sequence from $S$. Our objective is to construct a classification model which, given an amino acid sequence $s\in S$, assigns posterior probability that the sequence has the ability to carry out each particular function from $F$, or is involved in each particular biological process in $P$. Similarly, given a functional term $f$, or a biological process $p$, our objective is to find the most likely sequences associated with that function or process.

%%%%%%%%%%%%%%%%%%%%%%
%%%%%%%%%%%%%%%%%%%%%%
\subsection{Data sets}
\label{sec:fanngodata}
%%%%%%%%%%%%%%%%%%%%%%
%%%%%%%%%%%%%%%%%%%%%%

We used the Swiss-Prot database from May 2010 (v.15.15) \citep{bairoch2005universal}. A data set $D_{MFO}$ of $26,707$ protein sequences was generated by selecting sequences with molecular functions that were supported by at least one of the following experimental evidence codes: EXP, IDA, IPI, IMP, IGI, IEP, TAS, IC. The $26,707$ sequences in $D_{MFO}$ consisted of a total of $4,276$ molecular function terms associated with them. Similarly, a separate set $D_{BPO}$ consisting of $29,118$ sequences with a total of 11,300 associated biological process terms was generated using the same criterion. We note that $\left| D_{MFO}\cap D_{BPO} \right| = 19,240$ and $\left| D_{MFO} \cup D_{BPO} \right| = 36,585$.

To arrive at the final data sets, $1,429$ proteins with known function in Swiss-Prot v15.15 were removed. Such sequences were either shorter than $50$ amino acids or were associated with the same gene name as some other protein in the same organisms. Data sets are summarized in \Cref{tab:fanngoT1}.

%%%%%%%%%%%%%%%%%%%%%%
%%%%%%%%%%%%%%%%%%%%%%

\begin{table*}[t]
\caption{The number of proteins per species in each dataset.
\label{tab:fanngoT1}}
{
\begin{tabular}{lccc}
\hline
Organism & \hspace*{3cm} & MFO & BPO\\
\hline
\emph{H. sapiens} & \hphantom & $7,253$ & $6,515$\\
\emph{S. cerevisiae} & \hphantom & $4,062$ & $4,071$\\
\emph{M. musculus} & \hphantom & $3,945$ & $4,672$\\
\emph{R. norvegicus} & \hphantom & $2,624$ & $2,696$\\
\emph{A. thaliana} & \hphantom & $2,042$ & $2,813$\\
\emph{D. melanogaster} & \hphantom & $1,627$ & $1,591$\\
\emph{E. coli K-12} & \hphantom & $1,498$ & $1,009$\\
\emph{S. pombe} & \hphantom & $1,079$ & $1,935$\\
\emph{C. elegans}  & \hphantom& $569$ & $1,599$\\
All other & \hphantom & $2,008$ & $2,217$\\
\hline
Total & \hphantom & $26,707$ & $29,118$\\
\end{tabular}}{}

\end{table*}
%%%%%%%%%%%%%%%%%%%%%%
%%%%%%%%%%%%%%%%%%%%%%


%%%%%%%%%%%%%%%%%%%%%%
%%%%%%%%%%%%%%%%%%%%%%
\subsection{Data representation}
%%%%%%%%%%%%%%%%%%%%%%
%%%%%%%%%%%%%%%%%%%%%%

Sequence alignments were used to represent each protein sequence as a fixed-length vector in a feature space. Each dimension in the feature space was selected to correspond to one term from a set of available functional terms $F$ (or $P$ for the biological process ontology). While we tested several ways of encoding alignment data into features we found that using i-scores, as proposed by the GOtcha algorithm \citep{Martin2004}, as features worked the best. For the completeness of this work, we briefly summarize this representation. 

First, let $e(s, s_{i})$ be the E-value obtained by aligning target sequence $s$ to the $i$-th sequence in database of proteins with experimentally determined functions, i.e. $s_{i}\in S$. The r-score for functional term $f$ is then generated as 
\[
r_{f}(s)=-\underset{s_{i}\in S_{f}}{\sum}\log (e(s,s_{i}))+c
\]
where $S_{f}$ is a subset of $S$ containing all proteins with functional term $f$ and $c$ is a constant value added to the sum to ensure non-negative r-scores (here we used $c=2$ and E-value threshold of 10). The i-scores for each function were then calculated by normalizing $r_{f}$ by the r-score of the root node in a given ontology (term $f_{root}\in F$) as
\[
i_{f}=\frac{r_{f}}{r_{root}}
\]
Finally, a feature vector is obtained by concatenating i-scores for each of the $|F|$ functional terms. We note that the i-score for the root term always equals $1$; thus, it was excluded from the vector representation. For the molecular function ontology, the i-score vector representation consists of $|F|$ dimensions, while for the biological process ontology, the i-score representation consists of $|P|$ dimensions. These features were used for the basic version of Functional ANNotator (FANN-GO).

In order to take advantage of the fact that the quality of transferring functions from sequences within the same species was higher than that achieved when transferring functions only from sequences in different species, we made two additional sets of i-score features, one based on the i-scores using the proteins from the same species and another using the i-scores from the proteins from different species only. These features were used for the version of the predictor referred to FANN-GO$^{species}$. 

%%%%%%%%%%%%%%%%%%%%%%
%%%%%%%%%%%%%%%%%%%%%%
\subsection{Classification models}
%%%%%%%%%%%%%%%%%%%%%%
%%%%%%%%%%%%%%%%%%%%%%

In order to address the multi-label classification problem of protein function we utilized a multi-output feed-forward neural network framework. Multi-output networks have the ability to simultaneously learn multiple dependent target variables, a property that is well suited to the problem of predicting mutually non-exclusive terms of protein function ontologies. 

Before training a multi-output neural network we performed several data preprocessing steps. All features were first normalized using the z-score method. Feature selection filtering was then performed by using the t-test. Finally, principal component analysis was performed in order to combine highly correlated features (retained variance = $99\%$).

Owing to the high-memory requirement of a multi-output neural network, it was not practically possible to train a model with more than $1000$ outputs on a data set of size $|D_{MFO}|$ or $|D_{BPO}|$. To overcome this limitation, we created ensembles of $100$ networks such that in each network only 100 randomly selected outputs were considered. Prediction values for a test sequence for a particular function were finally calculated as an average over the output scores generated from the subset of networks that included the given term in their output layers. All neural networks had $100$ hidden neurons, employing the resilient propagation algorithm \citep{Riedmiller1993} in training (with at most $1,000$ epochs). The networks were implemented using MATLAB.

%%%%%%%%%%%%%%%%%%%%%%
%%%%%%%%%%%%%%%%%%%%%%
\subsection{Model selection and evaluation}
\label{sec:fanngomodels}
%%%%%%%%%%%%%%%%%%%%%%
%%%%%%%%%%%%%%%%%%%%%%

The accuracy of the model was estimated using 10-fold cross-validation. Parameter selection for each neural network was performed on a separate validation set, such that only the best performing parameter set was used on the test partition. Furthermore, separate sets of features were generated for each fold using only alignments with sequences in the training portion of the data. An individual BLAST database was built for each fold's set of training sequences in order to ensure that alignment E-values were not influenced by sequences in the test portion of the data.  All methods were evaluated by plotting precision-recall curves. FANN-GO was evaluated against three different strategies. The Global-SID and Local-SID strategies represent transfer by sequence similarity in which each prediction was generated by transferring functional terms directly from sequences with sequence identity to the query sequence greater than the threshold. In addition, the performance of FANN was compared to the GOtcha classifier \citep{Martin2004} as well as the Na\"ive classifier. The Na\"ive classifier predicts the terms according to their descending prior probabilities in the training data, i.e. the term occurring in $75\%$ of training sequences will be predicted with score $0.75$ for all target proteins.

Precision-recall curves were generated as follows. For each query sequence, a set of predictions over all $|F|$ functions was generated. A decision threshold $t$ value above which all predictions were taken was incrementally reduced from $1$ to $0$, in steps of $0.01$. Terms with predictions scores above a particular threshold $t_{i}$ were selected, and each term was propagated towards the root of the ontology. This resulted in a set of predicted terms $P$. The precision ($pr$) and recall ($rc$) between the predicted terms $P$ and true terms $T$ associated with sequence \emph{s} were then calculated as 

\[
pr=\frac{|T \cap P|}{| P |} 
\]
and 
\[
rc=\frac{| T\cap P |}{|T|}
\]
The final precision and recall were averaged over all test sequences to create a point in the precision-recall space.

